% !PRINTED 544
\weberpage{544}

If one carries out on the right-hand side of this formula the permutation $\pi$ in the indices of the $\alpha$, then by 1. there results
\[
\sum^{\varepsilon} \varepsilon^{-k-\nu}(\varepsilon,\alpha)^{\nu},
\]
that is, $A_k^{(\nu)}$ passes into $A_{k+\nu}^{(\nu)}$, as is asserted in Theorem 2.

Theorem 2. is a special case of a general theorem. If we develop a product of arbitrarily many factors
\[
(\varepsilon,\alpha)^{\nu}(\varepsilon^{\lambda_1},\alpha)^{\nu_1}(\varepsilon^{\lambda_2},\alpha)^{\nu_2}\dots,
\]
in which $\nu, \nu_1, \nu_2\dots$ are positive and $\lambda_1, \lambda_2\dots$ arbitrary integers, and order it, as before $(\varepsilon,\alpha)^{\nu}$, according to powers of $\varepsilon$, the result may be
\weq{7}{(\varepsilon,\alpha)^{\nu}(\varepsilon^{\lambda_1},\alpha)^{\nu_1}(\varepsilon^{\lambda_2},\alpha)^{\nu_2}\dots = \sum_{0,\,m-1}^{h}\varepsilon^h B_h,}

in which the $B_h$ are forms of the variables $\alpha$ independent of $\varepsilon$. Since this development also holds for all $m$th roots of unity $\varepsilon$, one can apply the formula (2) and obtains
% [sic] authority prints B_k rather than mB_k and puts each \lambda_j on \alpha in the
% factors below; coefficient extraction from (7) and the product displayed above require
% mB_k and (\varepsilon^{\lambda_j},\alpha), respectively.
\[
B_k = \sum^{\varepsilon} \varepsilon^{-k}(\varepsilon,\alpha)^{\nu}(\varepsilon,\alpha^{\lambda_1})^{\nu_1}(\varepsilon,\alpha^{\lambda_2})^{\nu_2}\dots
\]
from which one can conclude, by 1., that $B_k$ passes by the substitution $\pi$ into
\[
B_{k+\nu+\lambda_1\nu_1+\lambda_2\nu_2+\cdots}
\]
We have therefore:

\gesperrt{3. The permutation $\pi$, applied to the indices of the $\alpha$, calls forth among the indices of the coefficients of the product, ordered according to $\varepsilon$,}
\weq{8}{(\varepsilon,\alpha)^{\nu}(\varepsilon^{\lambda_1},\alpha)^{\nu_1}(\varepsilon^{2},\alpha)^{\nu_2}\dots} % [sic] authority prints \varepsilon^2 here, omitting the expected \lambda_2.
\gesperrt{the permutation $\pi^{\nu+\lambda_1\nu_1+\lambda_2\nu_2\cdots}$.}

In these theorems the permutation $\pi$ may be repeated; they remain correct when $\pi$ is replaced by any power of $\pi$, and they hold whatever $m$th root of unity $\varepsilon$ may be. But if $\varepsilon$ is an imprimitive root of unity, certain simplifications occur, with which we must still become acquainted. If $e$ is a divisor of $m$,
\[
m = ef
\]
and $\varepsilon$ is any $e$th root of unity, then the expression (1) becomes
\begin{align*}
(\varepsilon,\alpha) &= \alpha + \varepsilon\alpha_1 + \cdots + \varepsilon^{e-1}\alpha_{e-1}\\
&\quad + \alpha_e + \varepsilon\alpha_{e+1} + \cdots + \varepsilon^{e-1}\alpha_{2e-1}\\
&\quad \cdots\cdots\cdots\cdots\cdots\cdots\cdots\cdots\cdots\cdots\\
&\quad + \alpha_{e(f-1)} + \varepsilon\alpha_{e(f-1)+1} + \cdots + \varepsilon^{e-1}\alpha_{m-1}
\end{align*}
